\subsection{\texorpdfstring{Functions of class C \(^{r}\) and morphisms of manifolds}{Functions of class C\^{}\{r\} and morphisms of manifolds}}

5.3.1. Let X be a manifold of class \(\mathbf{C}^{r}\), let F be a separated polynormed space, and let \(f\) be a map from X into F. We say that \(f\) is of class \(\mathbf{C}^{r}\) if, for every chart \((V, \varphi, E)\) of X, the map \(f\circ\varphi^{-1}\) from \(\varphi(\mathrm V)\) into \(F\) is of class \(\mathbf{C}^{r}\). For this, it suffices that this condition be satisfied for the charts of an atlas of X. The set of maps of class \(\mathbf{C}^{r}\) from X into F forms a vector subspace of the space of all maps from X into F. We denote it by \(\mathcal{C}^{r}(\mathrm{X};\mathrm{F})\) .

When F = K, we set \(\mathcal{C}^{r}(X;K)=\mathcal{C}^{r}(X)\) ; it is a sub-K-algebra of the algebra of maps from X into K. The elements of \(\mathcal{C}^{r}(X)\) are also called morphic functions.

When X is an open subset of a Banach space, this terminology agrees with that of nos. 2.3.1, 3.2.1 and 4.2.1.

5.3.2. Let X and Y be two manifolds of class \(\mathbf{C}^{r}\) and let \(f\) be a map from X into Y. We say that \(f\) is of class \(\mathbf{C}^{r}\) or is a morphism of manifolds (of class \(\mathbf{C}^{r}\) ) if it is continuous and if, for every chart \((\mathbf{V},\psi ,\mathbf{F})\)

of Y, the map \(\psi \circ f\) from the open submanifold \(f^{-1}(\mathbf{V})\) into the Banach space F is of class \(\mathbf{C}^{r}\) . For this, it suffices that there exist an atlas \(\mathcal{A}\) of Y such that, for every chart \((\mathbf{V}, \psi, \mathbf{F}) \in \mathcal{A}\) , the set \(f^{-1}(\mathbf{V})\) is open in X and the map \(\psi \circ f\) from the open submanifold \(f^{-1}(\mathbf{V})\) of X into F is of class \(\mathbf{C}^{r}\) . The set of morphisms from X into Y is denoted by \(\mathcal{C}^{r}(\mathbf{X}; \mathbf{Y})\) . When Y is a Banach space equipped with its canonical manifold structure, the definitions of 5.3.1 and 5.3.2 are consistent. A map of class \(\mathbf{C}^{\omega}\) is also called a K-analytic map (or simply analytic). When \(\mathbf{K} = \mathbf{C}\) , we also say holomorphic map.

Let (U, \(\varphi\) , E) be a chart of X and (V, \(\psi\) , F) a chart of Y such that \(f(\mathbf{U}) \subset \mathbf{V}\) . The map \(\psi \circ f \circ \varphi^{-1}\) of the open subset \(\varphi(\mathbf{U})\) of E into the open subset \(\psi(\mathbf{V})\) of F is called the expression of \(f\) in the given charts.

5.3.3. Suppose X and Y are of finite dimension; let \(a \in \mathbf{X}\) , \(b \in \mathbf{Y}\) and \(f\) be a map from X into Y with \(f(a) = b\) . Suppose first that \(f\) is of class \(\mathbf{C}^r\) and consider coordinate systems \((\mathrm{U}; \xi^1, \ldots, \xi^m)\) of X at \(a\) and \((\mathrm{V}; \eta^1, \ldots, \eta^n)\) of Y at \(b\) respectively, with \(f(\mathrm{U}) \subset \mathrm{V}\) . Let \(\xi\) be the map \((\xi^1, \ldots, \xi^m)\) from U into \(\mathbf{K}^m\) . There then exist functions \(u^j\) of class \(\mathbf{C}^r\) on the open subset \(\xi(\mathrm{U})\) of \(\mathbf{K}^m\) , with values in K, such that the coordinates of a point \(y = f(x)\) of Y, (for \(x\) in U), are given by the formulas:

\[
\eta^ {j} (y) = u ^ {j} (\xi^ {1} (x), \dots , \xi^ {m} (x)) \quad \text { for } 1 \leqslant j \leqslant n\tag{1}
\]

which is equivalent to:

\[
\eta^ {j} \circ f = u ^ {j} (\xi^ {1}, \dots , \xi^ {m}) \quad \text { for } 1 \leqslant j \leqslant n.\tag{2}
\]

We say that the preceding formulas constitute the expression of f by means of the chosen coordinate systems.

Conversely, if, for every point \(a\) of X, one can find a coordinate system of X at \(a\) and a coordinate system of Y at \(b = f(a)\) satisfying the preceding conditions, then \(f\) is of class \(\mathbf{C}^{r}\) .

5.3.4. Morphisms of manifolds satisfy the axioms of Ens., chap. IV, § 2:

\begin{enumerate}
\def\labelenumi{\alph{enumi})}
\item
  The composite of two morphisms is a morphism.
\item
  For a bijection \(f:X\to Y\) to be an isomorphism, it is necessary and sufficient that \(f\) and \(f^{-1}\) are morphisms.
\end{enumerate}

5.3.5. Suppose that \(r = \omega\) . Let X be a manifold and \(f, g\) two analytic maps from X into a separated polynormed space or into a separated analytic manifold. The set of points in a neighborhood of which \(f\) and \(g\) are equal, is open and closed in X.

5.3.6. Suppose that K = R and \(r \neq \omega\) . A manifold X is said to admit partitions of unity of class \(C^{r}\) if, for every locally finite open covering of X, there exists a continuous partition of unity (Top. Gén., ch. IX, § 4, n° 3) subordinate to this covering and consisting of functions of class \(C^{r}\) .

Let E be a Banach space; consider the following property:

For any disjoint closed subsets A and B of E, there exists a function f of class \(C^{r}\) on E, with values in R, such that \(f(x) = 0\) for \(x \in A\) , \(f(x) = 1\) for \(x \in B\) and \(0 \leqslant f(x) \leqslant 1\) for all \(x \in E\) .

If S is a set of Banach spaces having property (PU), then every paracompact manifold of type S admits partitions of unity of class \(\mathbf{C}^{r}\).

Every finite-dimensional space, every separable Hilbert space, has property (PU).

5.3.7. Suppose K is ultrametric. Let X be a paracompact manifold. For every open cover U of X, there exists a partition of X into open and closed subsets subordinate to the cover U.

5.3.8. Suppose K is locally compact. Let X and Y be two pure manifolds of finite dimension, and let f be a morphism from X to Y. If dim X \textless{} dim Y and if the topology of X admits a countable base of open sets, \(f(X)\) is meager in Y.
% semantic-fragments: legacy-input-seam
\relax{}
